% !PRINTED 79
\weberpage{79}

Die so definirten Grössen
\[
        A_{\alpha_1,\alpha_2,\ldots,\alpha_\nu}^{\beta_1,\beta_2,\ldots,\beta_\nu}
\]
\gesperrt{heissen die \(\nu\)ten Unterdeterminanten oder Unterdeterminanten \(\nu\)ter Ordnung.} Sie sind dargestellt durch \((n-\nu)\)reihige Determinanten.

Aus III. folgt in Bezug auf diese Unterdeterminanten der Satz:

\gesperrt{IV. Die Unterdeterminante \(A_{\alpha_1,\alpha_2,\ldots,\alpha_\nu}^{\beta_1,\beta_2,\ldots,\beta_\nu}\) ändert nur ihr Vorzeichen, wenn zwei ihrer unteren oder zwei ihrer oberen Indices vertauscht werden, oder allgemeiner: sie bleibt dem absoluten Werthe nach ungeändert, wenn die Anordnung der Indices \(\alpha_1,\alpha_2\ldots\alpha_\nu\) durch irgend eine andere Anordnung ersetzt wird und ändert das Zeichen oder nicht, je nachdem diese Permutation zur zweiten oder zur ersten Art gehört.}

Bezeichnen wir aber mit \(\alpha_1',\alpha_2'\ldots\alpha_\nu'\) irgend eine Anordnung der \(\alpha_1,\alpha_2\ldots\alpha_\nu\), so enthält die Determinante \(A\) auch den Complex der Glieder
\[
        \pm a_{\alpha_1'}^{(\beta_1)}a_{\alpha_2'}^{(\beta_2)}\cdots
        a_{\alpha_\nu'}^{(\beta_\nu)}
        A_{\alpha_1,\alpha_2,\ldots,\alpha_\nu}^{\beta_1,\beta_2,\ldots,\beta_\nu},
\]
und wenn wir also alle diese Glieder sammeln, so erhalten wir den Complex:
\begin{equation}
        A_{\alpha_1,\alpha_2,\ldots,\alpha_\nu}^{\beta_1,\beta_2,\ldots,\beta_\nu}
        \;\Sigma\pm a_{\alpha_1}^{(\beta_1)}a_{\alpha_2}^{(\beta_2)}\cdots a_{\alpha_\nu}^{(\beta_\nu)}.
\tag{6}
\end{equation}
Die hier auftretende \(\nu\)-reihige Determinante
\[
        \Sigma\pm a_{\alpha_1}^{(\beta_1)}a_{\alpha_2}^{(\beta_2)}\cdots a_{\alpha_\nu}^{(\beta_\nu)}=
        \begin{vmatrix}
        a_{\alpha_1}^{(\beta_1)}&a_{\alpha_2}^{(\beta_1)}&\cdots&a_{\alpha_\nu}^{(\beta_1)}\\
        a_{\alpha_1}^{(\beta_2)}&a_{\alpha_2}^{(\beta_2)}&\cdots&a_{\alpha_\nu}^{(\beta_2)}\\
        \vdots&\vdots&&\vdots\\
        a_{\alpha_1}^{(\beta_\nu)}&a_{\alpha_2}^{(\beta_\nu)}&\cdots&a_{\alpha_\nu}^{(\beta_\nu)}
        \end{vmatrix}
\]
wollen wir die zu \(A_{\alpha_1,\alpha_2,\ldots,\alpha_\nu}^{\beta_1,\beta_2,\ldots,\beta_\nu}\) complementäre Unterdeterminante nennen und mit \(B_{\alpha_1,\alpha_2,\ldots,\alpha_\nu}^{\beta_1,\beta_2,\ldots,\beta_\nu}\) bezeichnen. Sie enthält genau die Zeilen und Colonnen, die in \(A_{\alpha_1,\alpha_2,\ldots,\alpha_\nu}^{\beta_1,\beta_2,\ldots,\beta_\nu}\) fehlen und
%CONT>>
